\subsection{Negation Statement}
\begin{definition}
    [Negation Statement]
    Given a statement R, the statement $\neg R$ is called the negation statement of R. 
\end{definition}

According to the De'Morgan Law, we already known the two negation statements. 
\begin{example}
    Given statements $P$ and $Q$,
    \[
        \neg (P \wedge Q) = \neg P \vee \neg Q
    \]
    \[
            \neg (P \vee Q) = \neg P \wedge \neg Q
    \]
\end{example}

In the following examples, we will discuss how to get the negation statement of a statement involved quantifiers. 
\begin{example}
    Given set X and element x, which x is an element of set X, and a statement P about element x. 
    \begin{equation}
        \neg (\forall x \in X, P(x)) \equiv \exists x \in X, \neg P(x) \label{eq2.8.1}
    \end{equation}
    \begin{equation}
        \neg (\exists x \in X, P(x)) \equiv \forall x \in X, \neg P(x) \label{eq2.8.2}
    \end{equation}
    

    If a statement involved multiple quantifiers, negate it involves serveral iterations of Equation (\ref{eq2.8.1}) and (\ref{eq2.8.2}).
    \begin{align*}
        \neg \left(\forall x \in \R, \exists y \in \R, y^3 = x\right) &\equiv \neg \exists x \in \R, \left(\exists y \in \R, y^3 = x\right) \\
        &\equiv \exists x \in \mathbb{R}, \forall y \in \mathbb{R}, \neg(y^3 = x) \\
        &\equiv \exists x \in \mathbb{R}, \forall y \in \mathbb{R}, y^3 \neq x
    \end{align*}
\end{example}

Let us discuss how to get the negation statement of a statement involved implication.
\begin{example}
    Given statement P and Q, according to the previous sectoin, we already known
    \[
        P \Rightarrow Q \equiv \neg (P \wedge \neg Q) \equiv \neg P \vee Q
    \]
    Therefore, according to De'Morgan Law
    \[
        \neg \left(P \Rightarrow Q\right) \equiv \neg (\neg P \vee Q) \equiv P \wedge \neg Q
    \]
\end{example}